
\subsubsection{Consistency-Based Uncertainty Quantification}

Consistency-based methods measure uncertainty by detecting generative inconsistency across multiple model samples. SelfCheckGPT \citep{manakul2023selfcheckgpt} pioneered this approach, using natural language inference and n-gram overlap to score consistency between an initial response and multiple resampled outputs. The method operates purely in the generation space without requiring external knowledge bases.

Subsequent work extended this paradigm. MIND Framework \citep{chen2024mind} analyzes internal language model states for real-time detection. MetaQA \citep{guo2024metaqa} applies metamorphic testing to hallucination detection. These methods share a common limitation: they provide uncertainty scores but lack statistical guarantees on calibration quality or coverage.

Consistency methods are computationally efficient, typically requiring 5--10 samples. However, they cannot answer questions requiring rigorous statistical bounds, limiting application in risk-sensitive domains.

\subsubsection{Statistical Uncertainty Quantification}

Statistical uncertainty quantification methods provide calibration guarantees through rigorous probabilistic frameworks. Conformal prediction (Vovk et al., 2005; Shafer \& Vovk, 2008) offers distribution-free coverage guarantees. Given miscoverage rate $\alpha$ and an exchangeable calibration set, conformal intervals contain the true label with probability at least 1-$\alpha$. COIN \citep{wang2025coin} applies conformal prediction to language model factuality, achieving coverage above 90\% with false discovery rate control.

FactTest \citep{nie2024facttest} extends conformal methods with hypothesis testing frameworks providing Type I and Type II error control. Conformal methods guarantee coverage independent of model architecture or training procedure. However, they require large labeled calibration sets (typically 500--1000 examples) and incur significant computational cost.

Statistical methods provide rigorous bounds but ignore epistemic structure. A query where the model is highly inconsistent receives the same statistical treatment as one where the model is certain, leading to computational inefficiency.

\subsubsection{Integrated Approaches}

Prior work has explored combining multiple uncertainty quantification signals, primarily through independent cascades. DiverseAgentEntropy \citep{liu2024diverseagent} estimates uncertainty via multi-agent voting. C-LoRA \citep{jin2024clora} introduces parameter-efficient contextual uncertainty through LoRA adapters. These methods apply different uncertainty quantification techniques sequentially or in parallel but lack joint calibration. Signals are computed independently and combined post-hoc through thresholding or weighted averaging.

The closest prior work is independent cascade, where consistency filtering precedes conformal calibration. However, cascades treat the two methods as separate stages. Consistency thresholds and conformal parameters are tuned independently on the same validation set, with no bidirectional information flow.

No prior work integrates consistency and statistical methods through joint calibration where each signal informs and updates the other.

\subsubsection{Theoretical Foundations}

The relationship between consistency-based and statistical uncertainty quantification methods has been unclear. Impossibility results \citep{karbasi2025impossibility} prove that absolute hallucination detection is computationally equivalent to language identification, requiring expert-labeled negative examples. This appears to contradict the empirical success of SelfCheckGPT-style methods.

This paradox may be resolved by recognizing that consistency methods measure epistemic structure, not absolute truth. The impossibility result applies to detecting hallucinations without any oracle. Consistency methods implicitly use the model itself as an imperfect oracle through multi-sample agreement. Conformal methods provide aleatoric bounds through calibration set statistics, not epistemic guarantees. The two methods may operate on complementary dimensions.

Surveys \citep{kang2025uq_survey} catalog uncertainty quantification methods but note the lack of unified theoretical frameworks. The present work addresses this gap by examining the relationship between epistemic and aleatoric uncertainty.
